% \documentclass[11pt]{article}

% %%%% CUSTOM PREAMBLE
% \usepackage{./preamble}

% %% ADD MISC DEFINITIONS HERE %%%%%%%%%%%%%%%%%%%%%%%%%%%%%%%%%%%

% %%%%%%%%%%%%%%%%%%%%%%%%%%%%%%%%%%%%%%%%%%%%%%%%%%%%%
% %\addtolength{\topmargin}{-1cm}
% %\addtolength{\textheight}{2cm}

% \begin{document}

%\maketitle

\newpage
 %%*****************************************************
\noindent
\textbf{CS4268 AY2025-26 Sem 2}
\\
\noindent
Lecturer: Divesh Aggarwal
\hfill
Lecture 10                      %%% FILL IN LECTURE NUMBER HERE
\hfill
Date: March 26, 2026          %%% FILL IN LECTURE DATE HERE


\noindent
\rule{\textwidth}{1pt}

\medskip

%%%%%%%%%%%%%%%%%%%%%%%%%%%%%%%%%%%%%%%%%%%%%%%%%%%%%%%%%%%%%%%%
%% BODY OF SCRIBE NOTES GOES HERE
%%%%%%%%%%%%%%%%%%%%%%%%%%%%%%%%%%%%%%%%%%%%%%%%%%%%%%%%%%%%%%%%

%%%Write lecture topic above.
\textbf{Instructions}\\
\begin{itemize}
    \item Try to leave the preamble unchanged, only make changes to \texttt{main.tex} and  \texttt{references.bib}. Add images to the \texttt{images} folder. 
    \item Share one (zipped) directory preserving the directory structure on Canvas. 
    \item Prepare organized notes by elaborating on the broad items provided in different sections below based on the lecture material completing any missing details from lectures.
    \item Feel free to modify structure and/or order of contents as you see fit; the below is just a guideline.
\end{itemize} 


\kk{Today is on search and grover. Add any details you deem helpful/interesting.}

\section{Grover's algorithm}
\begin{enumerate}
    \item Want to search for one of $M$ marked elements among $N$ elements. Given a Boolean function $f: [N]=\{0,1,\dots,N-1\} \longrightarrow \{0,1\}$. $f(x)=1$ if $x$ is one of those marked elements, otherwise $f(x)=0$.
    
    \item How is this done classically? For deterministic algorithms, worst case scenario requires $N$ queries. Probabilistically, if we randomly choose $k$ items to be queried, and require that we find (at least) a solution with probability as least $1-\varepsilon$, then we have (assuming $M/N \ll 1$) $1-(1-\frac{M}{N})^k \approx kM/N \geq 1-\varepsilon$, i.e. for a fixed error margin we need $k=\Omega(N/M)$ queries to solve the problem. 
    
    \item Grover's algorithm solves this problem using $\mathcal{O}(\sqrt{N/M})$ oracle queries, thus providing a \textit{quadratic speedup in query complexity}.

    \begin{algorithm}
    \caption{Grover's Algorithm}
    %\label{algo:TGD}
    \begin{algorithmic}[1]
    \Require Oracle $O_f$ to implement $f$ unitarily, where $f$ specifies the solutions.
    \State Initialize the state $\ket{0^n}\ket{0}$.
    \State Apply $H^{\otimes n}$ to the $n$ work qubits and $HX$ to the ancillary qubit:
    \[
        (H^{\otimes n} \otimes HX)\ket{0^n}\ket{0} = \frac{1}{\sqrt{2^n}}\sum_{x=0}^{2^n-1} \ket{x}\ket{-}.
    \]
    \State Apply the Grover operator $G=H^{\otimes n}(2\ket{0^n}\bra{0^n}-I)H^{\otimes n}O_f$. Iterate $k = \left\lceil \frac{\pi}{4}\sqrt{\frac{N}{M}} \right\rceil$ times:
    \[
        G^k \left( \frac{1}{\sqrt{2^n}}\sum_{x=0}^{2^n-1} \ket{x} \right) = \alpha \left( \frac{1}{\sqrt{N-M}} \sum_{x \in A} \ket{x} \right) + \beta \left( \frac{1}{\sqrt{M}} \sum_{x \in B} \ket{x} \right)
    \]
    where $|\alpha| \ll 1$, $|\beta| \approx 1$. (Ancillary qubit suppressed.)
    \State Measure the $n$ work qubits.
    \Ensure $x \in B$ (the solution set) with high probability.
    \end{algorithmic}
    \end{algorithm}

    \begin{figure}[h]
    \label{GroverCircuit}
        \centering
        \begin{quantikz}[wire types = {b},classical
        gap=0.1cm, column sep=0.3cm]
        \lstick{$\ket{0^n}$} && \gate{H^{\otimes n}} \slice{} & \gate[2]{G} & \gate[2]{G} & \midstick[2,brackets=none]{$\cdots$} & \gate[2]{G} & \gate[2]{G} \slice{} && \meter{}\\
        \lstick{\text{Ancilla}\quad$\ket{0}$} & \gate{X} & \gate{H} &&&&&&& \rstick{$\ket{-}$}
        \end{quantikz}
        \caption{Quantum circuit implementing Grover's algorithm.}
        %\label{fig:enter-label}
    \end{figure}

    \begin{figure}[h]
        \centering
        \begin{quantikz}[wire types = {b},classical
        gap=0.1cm, column sep=0.3cm]
        & \gate[2]{G} & \midstick[2,brackets=none]{=} & \gate[2]{O_f} & \gate{H^{\otimes n}} & \gate{2\ket{0^n}\bra{0^n}-I} & \gate{H^{\otimes n}} &\\ 
        &&&&&&&
        \end{quantikz}
        \caption{The Grover Operator.}
    \end{figure}

    \item Analysis of Grover:
    The operator $2\ket{0^n}\bra{0^n}-I$ is a conditional phase shift operator: upon its implementation every computational basis state except for $\ket{0^n}$ acquires a phase shift of $-1$. This can be implemented using the phase-kickback trick, where $f(0^n)=0$ and $f(x)=1$ for all $x \neq 0^n$. The Grover operator $G$ can thus be written as
    \[
        G = (2\ket{\psi}\bra{\psi}-I)O_f.
    \]
    After the preprocessing but before implementing $G$ (this stage is indicated by the dotted red line in Fig. \ref{GroverCircuit}), we have the initial state $\ket{\psi}$. Let us define the sets 
    \begin{align*}
        A = \{x \in \{0,1\}^n \;|\; f(x)=0\}\\
        B = \{x \in \{0,1\}^n \;|\; f(x)=1\}.
    \end{align*}
    (that is, $B$ is the set of solutions to the search problem) and their corresponding uniform superpositions
    \begin{align*}
        \ket{\alpha} &= \frac{1}{\sqrt{N-M}} \sum_{x \in A} \ket{x}\\
        \ket{\beta} &= \frac{1}{\sqrt{M}} \sum_{x \in B} \ket{x}.
    \end{align*}
    We can then rewrite
    \[
        \ket{\psi} = \sqrt{\frac{N-M}{N}}\ket{\alpha} + \sqrt{\frac{M}{N}}\ket{\beta}  
    \]
    as a vector in the 2D real subspace spanned by $\ket{\alpha}$ and $\ket{\beta}$. Note that $\braket{\beta}{\alpha}=0$. Now comes the crucial part of the analysis: Since $O_f$ marks the solutions by introducing a phase shift of $-1$, we have $O_f(a\ket{\alpha}+b\ket{\beta}) = a\ket{\alpha}-b\ket{\beta}$, i.e. $O_f = 2\ket{\alpha}\bra{\alpha}-1$ is also a reflection operator. Thus the Grover operator, being a composition of two reflection operators in $\mathbb{R}^2$ is a rotation operator (see lemma below).

    \begin{lemma}
        In $\mathbb{R}^2$, the composition of two reflections is a rotation.
    \end{lemma}
    \begin{proof}
        Any reflection about the subspace spanned by the unit vector $v$ is given by $R_v = 2vv^T-1$. One can easily check (by evaluating the inner product) that for any arbitrary unit vector $w$, we have $\theta_{v,w} = \theta_{v,R_v(w)} = \frac{1}{2}\theta_{w,R_v(w)}$. Note that $\text{det}(R_v) = -1$ (where we have used the fact that the eigenvalues of the projector $vv^T$ are $1$ and $0$), as is expected for a reflection.

        Given arbitrary unit vectors $v,w$ defining the lines of reflection, we now show that their product is a proper rotation, i.e. it is orthogonal and has determinant $1$. Both proofs are straightforward. First we have $(2vv^T-1)(2vv^T-1) = 1$, thus $(R_vR_w)^T (R_vR_w) = I$; second we have $\text{det}(AB) = \text{det}(A)\text{det}(B)$, so $\text{det}(R_vR_w) = 1$.
    \end{proof}

    \begin{proposition}
        In the $\ket{\alpha},\ket{\beta}$ basis, the Grover operator is represented as
        \[
            G = 
            \begin{bmatrix}
                \cos \theta & -\sin \theta\\
                \sin \theta &
                \cos \theta
            \end{bmatrix},
        \]
        where $\theta$ is such that $\cos(\theta/2) = \sqrt{\frac{N-M}{N}}$ and $\sin(\theta/2) = \sqrt{\frac{M}{N}}$.
    \end{proposition}
    \begin{proof}
        For ease of notation let us denote the coefficients of $\ket{\psi}$ as $a=\sqrt{\frac{N-M}{N}}$ and $b=\sqrt{\frac{M}{N}}$. Then
        \begin{align*}
            G\ket{\alpha} &= \left( 2(a^2\ket{\alpha}\bra{\alpha} + b^2\ket{\beta}\bra{\beta} + ab\ket{\alpha}\bra{\beta} + ab\ket{\beta}\bra{\alpha}) - I \right)O_f\ket{\alpha}\\
            &= (2a^2-1)\ket{\alpha} + (2ab)\ket{\beta}.
        \end{align*}
        The resulting coefficients remind us of the trigonometric formulae $\cos(2\theta) = 2\cos^2(\theta)-1$ and $\sin(2\theta) = 2\sin(\theta)\cos(\theta)$. Thus if we define $\theta$ such that $\cos(\theta/2) = a = \sqrt{\frac{N-M}{N}}$ and $\sin(\theta/2) = b = \sqrt{\frac{M}{N}}$, then $G\ket{\alpha} = \cos(\theta)\ket{\alpha} + \sin(\theta)\ket{\beta}$. Likewise, we can show that $G\ket{\beta} = -\sin(\theta)\ket{\alpha} + \cos(\theta)\ket{\beta}$.
    \end{proof}

    Starting from the state $\ket{\psi} = \cos\frac{\theta}{2}\ket{\alpha} + \sin\frac{\theta}{2}\ket{\beta}$ and applying $G$ gives us, after $k$ iterations,
    \[
        G^k\ket{\psi} = \cos\left( \frac{2k+1}{2}\theta\right)\ket{\alpha} + \sin\left( \frac{2k+1}{2}\theta\right)\ket{\beta}.
    \]
    To maximise the probability of obtaining an element in $B$ (i.e. getting a solution), we want that $\frac{2k+1}{2}\theta\ \approx \frac{\pi}{2}$. This leads us to choose $k = \lfloor \frac{\pi}{2\theta} - \frac{1}{2} \rceil$, where $\lfloor x \rceil$ denotes the rounding of $x$ to the nearest integer. Note that error is at most $\theta$, since if we are more than $\theta$ off from $\pi/2$ we would rotate further by another $\theta$. So smaller $M/N$ implies smaller $\theta$ implies more precision.

    \item It is also reasonable to assume $M \leq N/2$. In this situation, $\frac{\theta}{2} \geq \sin \frac{\theta}{2} = \sqrt{\frac{M}{N}}$, so
    \[
        k = \left\lfloor \frac{\pi}{2\theta} - \frac{1}{2} \right\rceil \leq \frac{\pi}{2\theta} \leq \left\lceil \frac{\pi}{4}\sqrt{\frac{N}{M}} \right\rceil,
    \]
    i.e. the number of Grover iterations required is $k = \mathcal{O}(\sqrt{\frac{N}{M}})$, a quadratic improvement over the $\mathcal{O}(\frac{N}{M})$ classical solution. If $M \geq N/2$, we simply double $N$ by introducing an ancilla qubit and extending $f$ s.t. the original marked items remain marked only when the ancilla is in the state $\ket{0}$.

    \item What have we done? We have effectively \textit{amplified} the probability of getting a solution, from $\sin^2\left(\frac{\theta}{2}\right)$ to $\sin^2\left(\frac{2k+1}{2}\theta \right) \approx 1$, by choosing $k$ appropriately. The full generality of this technique will be discussed further in the section on amplitude amplification.
\end{enumerate}













% \bibliographystyle{alpha}
% \bibliography{main.bib}
% \end{document}
